\documentclass[UTF8,11pt]{ctexart}
\usepackage[a4paper,margin=2.4cm]{geometry}
\usepackage{amsmath}
\usepackage{amssymb}
\usepackage{mathtools}
\usepackage{bm}
\usepackage{booktabs}
\usepackage{tabularx}
\usepackage{array}
\usepackage{float}
\usepackage{siunitx}
\usepackage{xcolor}
\usepackage{listings}
\usepackage{tikz}
\usepackage{hyperref}

\usetikzlibrary{arrows.meta,calc,positioning}
\hypersetup{hidelinks}
\sisetup{
      per-mode=symbol,
      group-separator={\,},
      exponent-product=\times,
      output-decimal-marker=.
}
\numberwithin{equation}{section}
\allowdisplaybreaks[2]
\setlength{\emergencystretch}{2em}
\renewcommand{\arraystretch}{1.22}

\definecolor{iceblue}{HTML}{DCEEFF}
\definecolor{waterblue}{HTML}{B9E3F5}
\definecolor{bathblue}{HTML}{EAF7FC}
\definecolor{fluxred}{HTML}{C73E1D}
\definecolor{gridgray}{HTML}{586069}
\definecolor{iceflowcodebg}{HTML}{F6F8FA}
\definecolor{iceflowcodeframe}{HTML}{D0D7DE}
\definecolor{iceflowcodekeyword}{HTML}{CF222E}
\definecolor{iceflowcodecomment}{HTML}{57606A}

\lstdefinestyle{iceflowcode}{
      language=Python,
      basicstyle=\ttfamily\small,
      keywordstyle=\color{iceflowcodekeyword}\bfseries,
      commentstyle=\color{iceflowcodecomment}\itshape,
      backgroundcolor=\color{iceflowcodebg},
      frame=single,
      rulecolor=\color{iceflowcodeframe},
      framerule=0.6pt,
      framesep=8pt,
      xleftmargin=0.02\linewidth,
      xrightmargin=0.02\linewidth,
      framexleftmargin=6pt,
      framexrightmargin=6pt,
      aboveskip=0.8em,
      belowskip=1em,
      columns=fullflexible,
      keepspaces=true,
      showstringspaces=false,
      breaklines=true,
      breakatwhitespace=false,
      tabsize=4
}

\newcommand{\fl}{f_{\ell}}
\newcommand{\dd}{\,\mathrm d}
\newcommand{\code}[1]{\texttt{#1}}

\title{固定冰二维融化算例的有限体积焓法实现\\
\large \texttt{stefan\_melting\_2d.py} 方法说明}
\author{}
\date{}

\begin{document}
\maketitle

\begin{abstract}
本文说明 IceFlow2D 中独立二维固定冰融化算例的物理模型、有限体积离散、焓--温度反演、
显式时间推进和验证方法。写作从单个网格单元的能量收支出发，不预设读者掌握有限体积法或
移动相界面算法。算例在恒温热水浴中央放置一块初始温度等于熔点的方冰，以体积焓作为唯一
守恒未知量；潜热由焓的等温平台表示，因此不必显式追踪冰--水边界。内部面热流采用调和平均
导热率，四壁采用半格距离恒温边界，所有单元由同一组面热流守恒更新。默认
$\SI{30}{\milli\metre}\times\SI{30}{\milli\metre}$、$120\times120$ 网格在
$\SI{180}{\second}$ 内得到清晰可见的二维融化，并以全局能量平衡、几何对称性和网格细化趋势
检查实现的自洽性。本文对应 \path{iceflow2d/stefan2d.py} 与
\path{iceflow2d/examples/stefan_melting_2d.py}，不涉及 LBM、流体速度或真实冰水密度差。
\end{abstract}

\noindent\textbf{关键词：} 二维融化；焓法；有限体积法；Stefan 相变；显式时间推进；能量守恒

\section{问题定义与阅读路线}

\subsection{求解对象}

计算域为二维矩形
\begin{equation}
      \Omega=[0,L_x]\times[0,L_y],
      \label{eq:domain}
\end{equation}
其中央初始冰区为
\begin{equation}
      \Omega_i^0=
      \left[\frac{L_x-w_i}{2},\frac{L_x+w_i}{2}\right]
      \times
      \left[\frac{L_y-h_i}{2},\frac{L_y+h_i}{2}\right].
      \label{eq:initial-ice-domain}
\end{equation}
初始冰温等于熔点 $T_m$，其余区域为温度 $T_b>T_m$ 的水；四面边界始终保持温度
$T_b$：
\begin{align}
      T(\bm x,0)&=
      \begin{cases}
            T_m, & \bm x\in\Omega_i^0,\\
            T_b, & \bm x\in\Omega\setminus\Omega_i^0,
      \end{cases}
      \label{eq:initial-temperature}\\
      T(\bm x,t)&=T_b,
      \qquad \bm x\in\partial\Omega.
      \label{eq:boundary-temperature}
\end{align}
图~\ref{fig:domain} 给出了几何和热边界。所谓“固定冰”只表示不求解冰块的平移、转动或受力；
冰区的相态并未固定，其边缘仍能吸热并逐渐融化。求解器也没有永久的“初始冰掩膜”：整个
冰--水区域都遵循同一焓闭合，因而热力学关系本身包含反向结冰分支。不过，当前
\code{Stefan2DConfig} 明确要求 $T_b>T_m$，且把初始水温和四壁温度统一设为 $T_b$，所以现有
命令行算例是专门的融化验证；若要构造结冰算例，还须分离初始水温与边界温度并相应放宽配置校验。

\begin{figure}[htbp]
      \centering
      \begin{tikzpicture}[>=Latex,scale=0.92]
            \fill[bathblue] (0,0) rectangle (8.8,6.4);
            \fill[iceblue] (2.9,2.0) rectangle (5.9,4.4);
            \draw[line width=1.0pt] (0,0) rectangle (8.8,6.4);
            \draw[line width=1.0pt,blue!65!black] (2.9,2.0) rectangle (5.9,4.4);
            \node at (4.4,3.2) {初始固定冰 $T=T_m$};
            \node at (1.25,5.25) {初始热水 $T=T_b$};
            \draw[<->] (0,-0.45) -- node[fill=white,inner sep=1pt] {$L_x$} (8.8,-0.45);
            \draw[<->] (-0.45,0) -- node[fill=white,inner sep=1pt,rotate=90] {$L_y$} (-0.45,6.4);
            \draw[<->] (2.9,1.65) -- node[fill=bathblue,inner sep=1pt] {$w_i$} (5.9,1.65);
            \draw[<->] (6.25,2.0) -- node[fill=bathblue,inner sep=1pt,rotate=90] {$h_i$} (6.25,4.4);
            \draw[fluxred,very thick,-{Latex[length=3mm]}] (4.4,6.95) -- (4.4,6.45);
            \draw[fluxred,very thick,-{Latex[length=3mm]}] (4.4,-0.95) -- (4.4,-0.45);
            \draw[fluxred,very thick,-{Latex[length=3mm]}] (-0.95,3.2) -- (-0.45,3.2);
            \draw[fluxred,very thick,-{Latex[length=3mm]}] (9.75,3.2) -- (9.25,3.2);
            \node[fluxred] at (4.4,7.25) {$T=T_b$};
            \node[fluxred] at (4.4,-1.25) {$T=T_b$};
            \node[fluxred,rotate=90] at (-1.25,3.2) {$T=T_b$};
            \node[fluxred,rotate=90] at (10.05,3.2) {$T=T_b$};
            \draw[->] (0.25,0.25) -- (1.2,0.25) node[right] {$x$};
            \draw[->] (0.25,0.25) -- (0.25,1.2) node[above] {$y$};
      \end{tikzpicture}
      \caption{二维固定冰融化算例。红色箭头只表示四壁向域内供热，不表示流体运动。}
      \label{fig:domain}
\end{figure}

\subsection{三种状态量及其职责}

初学时最重要的是区分以下三个量：
\begin{itemize}
      \item 温度 $T$ 回答“单元有多冷或多热”；
      \item 液相率 $\fl\in[0,1]$ 回答“单元内有多少比例已经成为水”，其中
            $\fl=0$ 为纯冰，$\fl=1$ 为纯水；
      \item 体积焓 $H$ 把显热和潜热合并成一个守恒能量变量，单位为
            $\si{\joule\per\metre\cubed}$。
\end{itemize}
程序每一步只沿着
\begin{equation}
      H^n\longrightarrow(T^n,\fl^n)
      \longrightarrow k^n
      \longrightarrow \bm q_f^n
      \longrightarrow H^{n+1}
      \label{eq:data-flow}
\end{equation}
推进。换言之，先由已有能量判断温度和相态，再算各面的热流，最后依据净流入热量更新能量。

\subsection{建模假设与适用边界}

本基准有意作如下简化：
\begin{enumerate}
      \item 仅考虑二维热传导，不求解速度、压力、自然对流或表面张力；
      \item 冰和水采用共同密度 $\rho$，默认值为 $\SI{1000}{\kilogram\per\metre\cubed}$；
      \item 冰为纯物质，熔化发生在单一温度 $T_m$，不设置人为的有限温宽糊状区；
      \item 按单位出平面厚度积分，故二维总能量的单位为
            $\si{\joule\per\metre}$；
      \item 材料参数在每个纯相内为常数，且不考虑辐射、黏性耗散和体积热源。
\end{enumerate}
因此，本例的定位是“二维焓离散的守恒、对称性与网格细化自洽验证”，而不是含流动的融冰
预测，也不是与一维解析 Stefan 相似解逐点比较。等密度假设还意味着它不能验证真实冰水密度差
引起的体积收缩、自由液面改变或质量交换。

\section{从能量收支到焓法}

\subsection{一个区域为什么会升温或融化}

设 $C\subset\Omega$ 为任意固定区域，$\bm n$ 为其外法向量，傅里叶热流为
\begin{equation}
      \bm q=-k\nabla T,
      \label{eq:fourier-law}
\end{equation}
其中 $k$ 是导热率，$\bm q$ 的单位为 $\si{\watt\per\metre\squared}$。负号表示热量从高温处
流向低温处。能量守恒写为
\begin{equation}
      \frac{\dd}{\dd t}\int_C H\dd A
      =-\oint_{\partial C}\bm q\cdot\bm n\dd s.
      \label{eq:integral-energy}
\end{equation}
右端是穿过边界流入区域的净热功率。因为采用单位出平面厚度，式~\eqref{eq:integral-energy}
两端单位均为 $\si{\watt\per\metre}$。下面逐步把该积分守恒式转化为单点处的微分方程。

首先，$C$ 是位置和形状均不随时间改变的固定控制体，因此时间导数可以移入面积积分：
\begin{equation}
      \frac{\dd}{\dd t}\int_C H\dd A
      =\int_C\frac{\partial H}{\partial t}\dd A.
      \label{eq:fixed-control-time-derivative}
\end{equation}
其次，对足够光滑的热流场 $\bm q=(q_x,q_y)$ 应用二维散度定理：
\begin{equation}
      \oint_{\partial C}\bm q\cdot\bm n\dd s
      =\int_C\nabla\cdot\bm q\dd A
      =\int_C\left(
      \frac{\partial q_x}{\partial x}
      +\frac{\partial q_y}{\partial y}
      \right)\dd A.
      \label{eq:divergence-theorem}
\end{equation}
式~\eqref{eq:divergence-theorem} 表明，边界上的净向外热流等于区域内部热流散度的面积积分。
将式~\eqref{eq:fixed-control-time-derivative} 和式~\eqref{eq:divergence-theorem} 代入
式~\eqref{eq:integral-energy}，依次得到
\begin{align}
      \int_C\frac{\partial H}{\partial t}\dd A
      &=-\int_C\nabla\cdot\bm q\dd A,
      \label{eq:integral-energy-after-divergence}\\
      \int_C\left(
      \frac{\partial H}{\partial t}
      +\nabla\cdot\bm q
      \right)\dd A
      &=0.
      \label{eq:zero-integral-energy}
\end{align}
控制体 $C$ 可以任意选取。若括号内的连续函数在某一点不为零，总能在该点附近选取一个足够小
的控制体，使积分不为零；这与式~\eqref{eq:zero-integral-energy} 矛盾。因此括号内的量必须
处处为零，即得到局部能量守恒方程
\begin{equation}
      \frac{\partial H}{\partial t}
      +\nabla\cdot\bm q=0.
      \label{eq:local-energy-conservation}
\end{equation}
最后代入傅里叶定律 $\bm q=-k\nabla T$：
\begin{align}
      \frac{\partial H}{\partial t}
      +\nabla\cdot(-k\nabla T)&=0,
      \label{eq:substitute-fourier-law}\\
      \frac{\partial H}{\partial t}
      &=\nabla\cdot(k\nabla T).
      \label{eq:enthalpy-pde}
\end{align}
把最后一式在二维直角坐标中完全展开，可写成
\begin{equation}
      \frac{\partial H}{\partial t}
      =
      \frac{\partial}{\partial x}
      \left(k\frac{\partial T}{\partial x}\right)
      +
      \frac{\partial}{\partial y}
      \left(k\frac{\partial T}{\partial y}\right).
      \label{eq:enthalpy-pde-2d}
\end{equation}

若材料不发生相变，输入热量通常使温度升高。这部分能够通过温度变化直接观察到的能量称为
显热；在比热近似为常数时，
\begin{equation}
      Q_s=mc_p\Delta T,
      \label{eq:sensible-heat}
\end{equation}
其中 $c_p$ 是定压比热，$\Delta T$ 是温度变化。

冰已经到达熔点 $T_m$ 后，继续输入的热量不再立即提高温度，而是用于克服固相结构的束缚，
使冰在相同温度下转变为水。这部分在相变过程中不能由温度升高直接表现出来、但仍真实储存在
材料焓中的能量称为\textbf{潜热}；对冰融化而言，更具体地称为\textbf{熔化潜热}。若质量
$m$ 的冰从“熔点处的纯冰”完全变成“熔点处的纯水”，所需潜热为
\begin{equation}
      Q_L=mL,
      \label{eq:latent-heat}
\end{equation}
其中 $L$ 是单位质量熔化潜热，单位为 $\si{\joule\per\kilogram}$；$m$ 的单位为
$\si{\kilogram}$，所以 $Q_L$ 的单位为 $\si{\joule}$。默认
$L=\SI{334000}{\joule\per\kilogram}$，表示 $\SI{1}{\kilogram}$、温度恰为
$\SI{0}{\degreeCelsius}$ 的冰，在温度仍保持 $\SI{0}{\degreeCelsius}$ 的条件下完全熔化，
需要吸收 $\SI{334}{\kilo\joule}$ 热量。这里没有包含先把低温冰加热至熔点或继续把融水升温
所需的显热。

若体积为 $V$ 的单元只融化一部分，液相率增加 $\Delta\fl$，则发生相变的质量和所吸收的
潜热分别为
\begin{align}
      \Delta m&=\rho V\Delta\fl,\\
      \Delta Q_L&=\Delta mL
      =\rho VL\Delta\fl.
      \label{eq:partial-latent-heat}
\end{align}
除以单元体积 $V$，可得到相应的体积焓增量
\begin{equation}
      \Delta H_L=\frac{\Delta Q_L}{V}
      =\rho L\Delta\fl.
      \label{eq:volumetric-latent-heat}
\end{equation}
这正是后文相变平台 $H=\rho L\fl$ 的来源。反向结冰时，相同质量的水会释放同样大小的潜热：
对材料本身而言焓减少，对周围环境而言则获得热量。因此，这种“能量改变而温度暂时不变”的
状态不能只用温度唯一表示，程序必须选取焓 $H$ 作为守恒未知量。

\subsection{焓、温度和液相率的闭合关系}

以熔点处的固态冰为零焓参考，并记冰、水定压比热为 $c_{p,i}$、$c_{p,w}$，则代码采用
\begin{equation}
      H(T,\fl)=
      \begin{cases}
            \rho c_{p,i}(T-T_m),
            & T<T_m,\quad \fl=0,\\[2mm]
            \rho L\fl,
            & T=T_m,\quad 0\le\fl\le1,\\[2mm]
            \rho L+\rho c_{p,w}(T-T_m),
            & T>T_m,\quad \fl=1.
      \end{cases}
      \label{eq:enthalpy-definition}
\end{equation}
式中 $H=0$ 只是所选能量零点，并不表示材料“没有能量”。相变平台的宽度为
$\rho L$：当 $H$ 从 $0$ 增加至 $\rho L$ 时，温度保持 $T_m$，液相率从 $0$ 增加至 $1$。
图~\ref{fig:enthalpy-curve} 以几何方式展示了这一关系。

\begin{figure}[htbp]
      \centering
      \begin{tikzpicture}[x=1.0cm,y=0.82cm,>=Latex]
            \draw[->] (-4.0,0) -- (4.8,0) node[right] {$T$};
            \draw[->] (0,-2.8) -- (0,6.5) node[above] {$H$};
            \draw[dashed,gray] (0,-2.4) -- (0,6.0);
            \draw[dashed,gray] (-3.6,4.0) -- (4.4,4.0);
            \draw[very thick,blue!70!black] (-3.2,-2.3) -- (0,0);
            \draw[very thick,blue!70!black] (0,0) -- (0,4.0);
            \draw[very thick,blue!70!black] (0,4.0) -- (4.0,6.0);
            \fill[blue!70!black] (0,0) circle (1.7pt);
            \fill[blue!70!black] (0,4.0) circle (1.7pt);
            \node[below right] at (0,0) {$T_m$};
            \node[left] at (0,4.0) {$\rho L$};
            \node[rotate=34] at (-2.2,-1.2) {纯冰显热};
            \node[rotate=90,fill=white,inner sep=1pt] at (0.35,2.0) {等温相变};
            \node[rotate=27] at (2.7,5.25) {纯水显热};
            \draw[<->,fluxred] (0.75,0.15) -- node[right] {$\rho L$} (0.75,3.85);
            \node[align=left,anchor=west] at (1.15,2.0)
            {$0<\fl<1$\\表示单元平均相态};
      \end{tikzpicture}
      \caption{以熔点固冰为零点的焓--温度关系。竖直段不是有限温宽材料层，而是纯物质的潜热平台。}
      \label{fig:enthalpy-curve}
\end{figure}

数值程序已知 $H$ 后，需要反向恢复 $T$ 和 $\fl$。式~\eqref{eq:enthalpy-definition} 的逐单元
反演为
\begin{equation}
      (T,\fl)=
      \begin{cases}
            \left(T_m+\dfrac{H}{\rho c_{p,i}},0\right),
            & H<0,\\[3mm]
            \left(T_m,\dfrac{H}{\rho L}\right),
            & 0\le H\le\rho L,\\[3mm]
            \left(T_m+\dfrac{H-\rho L}{\rho c_{p,w}},1\right),
            & H>\rho L.
      \end{cases}
      \label{eq:enthalpy-inversion}
\end{equation}
有限网格中出现 $0<\fl<1$，表示一个控制体被冰--水界面部分占据后的体积平均，不表示实际
材料在一个人为温度区间内逐渐软化。单元导热率则在两纯相值之间线性插值：
\begin{equation}
      k(\fl)=(1-\fl)k_i+\fl k_w.
      \label{eq:cell-conductivity}
\end{equation}
式~\eqref{eq:cell-conductivity} 是“由相态得到单元性质”；后文内部面上的调和平均是“由两个
单元性质得到面性质”，二者承担不同任务，不应混为同一种平均。

\section{从控制体能量平衡到二维离散}

\subsection{网格与单元平均量}

将计算域划分为 $N_x\times N_y$ 个矩形控制体，网格尺度和单元中心为
\begin{align}
      \Delta x&=\frac{L_x}{N_x},
      &x_i&=\left(i+\frac12\right)\Delta x,
      &&i=0,\ldots,N_x-1,\\
      \Delta y&=\frac{L_y}{N_y},
      &y_j&=\left(j+\frac12\right)\Delta y,
      &&j=0,\ldots,N_y-1.
      \label{eq:grid}
\end{align}
代码数组采用 $[j,i]$ 顺序，其形状为 $(N_y,N_x)$。$H_{ij}$ 表示单元内体积焓的平均值；
在单位出平面厚度下，该单元储存的相对能量为
\begin{equation}
      E_{ij}=H_{ij}\Delta x\Delta y,
      \qquad [E_{ij}]=\si{\joule\per\metre}.
      \label{eq:cell-energy}
\end{equation}

有限体积法的核心不是在单元中心直接近似式~\eqref{eq:enthalpy-pde} 的所有导数，而是把
式~\eqref{eq:integral-energy} 应用于每一个小矩形：分别算四个面送入或带走多少热量，再求净值。
图~\ref{fig:control-volume} 中，$P=(i,j)$ 是当前单元，$E,W,N,S$ 为相邻单元；面热流的正方向
统一取坐标正方向。

\begin{figure}[htbp]
      \centering
      \begin{tikzpicture}[>=Latex,scale=1.05]
            \foreach \x in {0,2,4} {\draw[gridgray] (\x,0) -- (\x,6);}
            \foreach \y in {0,2,4,6} {\draw[gridgray] (0,\y) -- (6,\y);}
            \fill[iceblue,opacity=0.75] (2,2) rectangle (4,4);
            \node at (3,3) {$P=(i,j)$};
            \node at (1,3) {$W$};
            \node at (5,3) {$E$};
            \node at (3,5) {$N$};
            \node at (3,1) {$S$};
            \draw[fluxred,very thick,-{Latex[length=3mm]}] (2.25,3.55) -- (1.25,3.55)
                  node[above,midway] {$q^x_{i-\frac12,j}$};
            \draw[fluxred,very thick,-{Latex[length=3mm]}] (3.75,2.45) -- (4.75,2.45)
                  node[below,midway] {$q^x_{i+\frac12,j}$};
            \draw[fluxred,very thick,-{Latex[length=3mm]}] (2.55,2.25) -- (2.55,1.25)
                  node[left,midway] {$q^y_{i,j-\frac12}$};
            \draw[fluxred,very thick,-{Latex[length=3mm]}] (3.45,3.75) -- (3.45,4.75)
                  node[right,midway] {$q^y_{i,j+\frac12}$};
            \draw[<->] (2,6.35) -- node[fill=white,inner sep=1pt] {$\Delta x$} (4,6.35);
            \draw[<->] (6.35,2) -- node[fill=white,inner sep=1pt,rotate=90] {$\Delta y$} (6.35,4);
            \node[align=left,anchor=west] at (7.0,4.7)
            {正号约定：\\$q^x>0$ 指向 $+x$，\\$q^y>0$ 指向 $+y$};
      \end{tikzpicture}
      \caption{二维单元中心有限体积及其四个面热流。箭头给出代码中的坐标正方向，并不表示所有热流实际为正。}
      \label{fig:control-volume}
\end{figure}

\subsection{内部面热流与调和平均}

以 $P=(i,j)$ 和东侧 $E=(i+1,j)$ 之间的面为例，中心温度梯度用两点差分近似：
\begin{equation}
      q^x_{i+\frac12,j}
      =-k_{i+\frac12,j}
      \frac{T_{i+1,j}-T_{ij}}{\Delta x}.
      \label{eq:east-flux}
\end{equation}
若两侧导热率不同，热量依次穿过 $P$ 内的半格和 $E$ 内的半格。这两个热阻串联：
\begin{equation}
      R_{i+\frac12,j}
      =\frac{\Delta x/2}{k_{ij}}
      +\frac{\Delta x/2}{k_{i+1,j}}.
      \label{eq:series-resistance}
\end{equation}
令单一等效材料在距离 $\Delta x$ 上产生相同热阻，即
$R_{i+\frac12,j}=\Delta x/k_{i+\frac12,j}$，可得调和平均
\begin{equation}
      k_{i+\frac12,j}
      =\frac{2k_{ij}k_{i+1,j}}{k_{ij}+k_{i+1,j}}.
      \label{eq:harmonic-x}
\end{equation}
这一定义能正确反映“低导热一侧控制串联传热”的物理事实。$y$ 方向面完全同理：
\begin{align}
      k_{i,j+\frac12}
      &=\frac{2k_{ij}k_{i,j+1}}{k_{ij}+k_{i,j+1}},\\
      q^y_{i,j+\frac12}
      &=-k_{i,j+\frac12}
      \frac{T_{i,j+1}-T_{ij}}{\Delta y}.
      \label{eq:harmonic-y}
\end{align}

\subsection{四面恒温边界为何出现系数 2}

边界温度给在几何壁面，而第一层单元温度给在距壁面半格的中心。以左壁为例，传热距离是
$\Delta x/2$，故
\begin{equation}
      q^x_{-\frac12,j}
      =-k_w\frac{T_{0j}-T_b}{\Delta x/2}
      =\frac{2k_w}{\Delta x}(T_b-T_{0j}).
      \label{eq:left-boundary-flux}
\end{equation}
其余三壁按相同正方向约定写为
\begin{align}
      q^x_{N_x-\frac12,j}
      &=\frac{2k_w}{\Delta x}(T_{N_x-1,j}-T_b),
      \label{eq:right-boundary-flux}\\
      q^y_{i,-\frac12}
      &=\frac{2k_w}{\Delta y}(T_b-T_{i0}),
      \label{eq:bottom-boundary-flux}\\
      q^y_{i,N_y-\frac12}
      &=\frac{2k_w}{\Delta y}(T_{i,N_y-1}-T_b).
      \label{eq:top-boundary-flux}
\end{align}
这里的系数 $2$ 纯粹来自中心到壁面的半格距离。代码在边界半格固定使用水的导热率 $k_w$，
表示外壁与恒温水浴接触；它不是面平均的一般公式。默认工况的外层单元保持液态，因此该选择与
场景一致，但若要模拟固体直接接触壁面，则应重新定义边界热阻。

\subsection{守恒的显式单元更新}

在时间区间 $[t^n,t^{n+1}]$ 内把各面热流暂视为时刻 $n$ 的常数，式
~\eqref{eq:integral-energy} 给出
\begin{equation}
\boxed{
      H_{ij}^{n+1}=H_{ij}^{n}+\Delta t_n
      \left[
      \frac{q^x_{i-\frac12,j}-q^x_{i+\frac12,j}}{\Delta x}
      +
      \frac{q^y_{i,j-\frac12}-q^y_{i,j+\frac12}}{\Delta y}
      \right].
}
      \label{eq:explicit-update}
\end{equation}
例如，东面为正的 $q^x_{i+1/2,j}$ 把热量从当前单元带走，因而在式
~\eqref{eq:explicit-update} 中带负号。同一内部面通量只计算一次：它是左单元的流出，同时是右单元
的流入，所以全域求和时逐面严格抵消。该“共享面通量”正是有限体积离散能保持全局能量守恒的
关键。

\section{时间推进、稳定性与完整算法}

\subsection{显式步长限制}

显式方法直接由旧时刻状态计算新时刻状态，不需求解联立方程，但时间步不能过大。程序先构造
覆盖纯冰、纯水及插值单元的保守扩散率上界
\begin{equation}
      \alpha_b
      =\frac{\max(k_i,k_w)}
      {\rho\min(c_{p,i},c_{p,w})}.
      \label{eq:diffusivity-bound}
\end{equation}
注意式~\eqref{eq:diffusivity-bound} 是便于稳定性判断的上界；对任意材料参数，它不一定恰好等于
两个纯相真实热扩散率中的最大值。考虑半格 Dirichlet 边界后，程序采用
\begin{equation}
      \Delta t_{\max}
      =\frac{1}
      {3\alpha_b\left(\Delta x^{-2}+\Delta y^{-2}\right)}.
      \label{eq:maximum-time-step}
\end{equation}
默认候选步长由 Fourier 数控制：
\begin{equation}
      \Delta t_0
      =\mathrm{Fo}\,
      \frac{\min(\Delta x,\Delta y)^2}{\alpha_b},
      \qquad 0<\mathrm{Fo}\le\frac16.
      \label{eq:default-time-step}
\end{equation}
对方形网格，式~\eqref{eq:maximum-time-step} 化为
$\Delta t_{\max}=\Delta x^2/(6\alpha_b)$，恰对应 $\mathrm{Fo}\le1/6$。默认
$\mathrm{Fo}=0.15$ 留有一定裕量。这个限制是数值稳定条件，不代表真实融化过程存在同样的最小
物理时间尺度；减小步长主要改变离散误差与计算成本。

\subsection{准确命中输出时刻}

配置步长 $\Delta t_0$ 通常不能整除 $\SI{30}{\second}$ 的输出间隔。若下一输出时刻为
$t_{\mathrm{out}}$，实际使用
\begin{equation}
      \Delta t_n=min(\Delta t_0,t_{\mathrm{out}}-t_n).
      \label{eq:short-step}
\end{equation}
因此每段末尾可能插入一个较短步，快照仍精确落在指定物理时刻。这也解释了默认总步数不能简单
写成 $\lceil t_{\mathrm{end}}/\Delta t_0\rceil$。

\subsection{求解流程}

实现中没有显式界面追踪、非线性迭代或隐式矩阵求解。完整算法可概括为：
\begin{enumerate}
      \item 校验物理参数、网格和冰区对齐关系；建立单元中心坐标与初始冰掩膜；
      \item 由初始温度和相态构造 $H^0$，再用式~\eqref{eq:enthalpy-inversion} 恢复
            $T^0,\fl^0$；
      \item 对每一时间步，用式~\eqref{eq:cell-conductivity} 计算单元导热率；
      \item 用式~\eqref{eq:harmonic-x}--\eqref{eq:harmonic-y} 计算内部面热流，用式
            ~\eqref{eq:left-boundary-flux}--\eqref{eq:top-boundary-flux} 计算四壁热流；
      \item 用式~\eqref{eq:explicit-update} 一次性更新全域焓，同时由同一组边界热流累计输入热量；
      \item 由新焓反演 $T^{n+1},\fl^{n+1}$，检查所有值有限且 $\fl\in[0,1]$；
      \item 到达输出时刻后记录场量、能量、固相面积与对称性诊断，直至终止时间。
\end{enumerate}

下列伪代码与 \code{Stefan2DSolver.step()} 的调用顺序一致：
\begin{lstlisting}[style=iceflowcode]
while time < next_output_time:
    dt = min(configured_dt, next_output_time - time)
    k_cell = (1.0 - liquid_fraction) * k_ice \
             + liquid_fraction * k_water
    flux_x, flux_y = conductive_face_fluxes(T, k_cell, T_bath)
    H += dt * ((flux_x[:, :-1] - flux_x[:, 1:]) / dx
               + (flux_y[:-1, :] - flux_y[1:, :]) / dy)
    boundary_heat += dt * boundary_power(flux_x, flux_y)
    T, liquid_fraction = recover_state_from_enthalpy(H)
    time += dt
record_snapshot()
\end{lstlisting}

\section{初始化、几何对齐与代码映射}

\subsection{几何栅格化与初始焓}

初始冰掩膜在单元中心按
\begin{equation}
      \left|x_i-\frac{L_x}{2}\right|<\frac{w_i}{2},
      \qquad
      \left|y_j-\frac{L_y}{2}\right|<\frac{h_i}{2}
      \label{eq:ice-mask}
\end{equation}
生成。配置预先要求冰的四条边落在网格面上，因此式~\eqref{eq:ice-mask} 的严格不等号不会产生
边界单元归属歧义。设
\begin{equation}
      N_{i,x}=\frac{w_i}{\Delta x},\qquad
      N_{i,y}=\frac{h_i}{\Delta y},
      \label{eq:ice-cell-count}
\end{equation}
则 $N_{i,x},N_{i,y}$ 必须为至少 $2$ 的整数，且居中边距格数
\begin{equation}
      \frac{N_x-N_{i,x}}{2},\qquad
      \frac{N_y-N_{i,y}}{2}
      \label{eq:margin-cell-count}
\end{equation}
也必须为整数并不小于 $1$。默认方冰覆盖 $48\times48$ 个单元，每侧有 $36$ 个水浴单元。
初始化结束后还核对离散固相面积是否等于 $w_i h_i$。

初始冰恰处于熔点且尚未融化，所以其焓为零；初始水是温度为 $T_b$ 的完全液态水：
\begin{equation}
      H_{ij}^{0}=
      \begin{cases}
            0,
            & (x_i,y_j)\in\Omega_i^0,\\[1mm]
            \rho\left[L+c_{p,w}(T_b-T_m)\right],
            & (x_i,y_j)\in\Omega\setminus\Omega_i^0.
      \end{cases}
      \label{eq:initial-enthalpy}
\end{equation}
式~\eqref{eq:initial-enthalpy} 中水的焓包含先熔化所需的参考潜热 $\rho L$，再加上从 $T_m$
升至 $T_b$ 的显热。

\subsection{配置防错条件}

\code{Stefan2DConfig} 在计算开始前检查：网格数是至少 $4$ 的整数；域尺寸、冰尺寸、材料参数、
终止时间和输出间隔为有限正数；$T_b>T_m$；冰严格小于计算域；式
~\eqref{eq:ice-cell-count}--\eqref{eq:margin-cell-count} 满足对齐条件；$0<\mathrm{Fo}\le1/6$；
用户指定的 $\Delta t$ 不超过式~\eqref{eq:maximum-time-step}。这些检查的作用不是改变物理模型，
而是在计算前拒绝可能导致模糊几何或显式发散的输入。

\subsection{主要程序对象}

表~\ref{tab:code-map} 给出数学量与代码对象的对应关系。由于所有字段均为 NumPy 二维数组，
本独立例子不调用 Taichi 或 LBM。

\begin{table}[htbp]
      \centering
      \caption{数学量与主要代码对象的映射。}
      \label{tab:code-map}
      \begin{tabularx}{\textwidth}{>{\raggedright\arraybackslash}p{0.20\textwidth}
            >{\raggedright\arraybackslash}p{0.33\textwidth}X}
            \toprule
            数学对象 & 代码对象 & 作用 \\
            \midrule
            物理与数值参数 & \code{Stefan2DConfig} & 校验几何、材料、输出时间和稳定步长 \\
            $H_{ij}$ & \code{enthalpy\_j\_m3} & 唯一守恒热力学状态，形状为 $(N_y,N_x)$ \\
            $T_{ij}$ & \code{temperature\_c} & 由焓分段反演得到的单元中心温度 \\
            $f_{\ell,ij}$ & \code{liquid\_fraction} & 由焓反演得到的液相体积分数 \\
            $q^x$ & \code{\_flux\_x\_w\_m2} & 所有竖直面的热流，形状为 $(N_y,N_x+1)$ \\
            $q^y$ & \code{\_flux\_y\_w\_m2} & 所有水平面的热流，形状为 $(N_y+1,N_x)$ \\
            单步更新 & \code{Stefan2DSolver.step()} & 面热流、焓更新、边界热量累计和状态反演 \\
            全程推进 & \code{Stefan2DSolver.run()} & 构造精确输出时刻并保存快照 \\
            文件与绘图 & \path{examples/stefan_melting_2d.py} & 参数解析、CSV/NPZ/JSON 与汇总图输出 \\
            \bottomrule
      \end{tabularx}
\end{table}

\section{守恒、融化程度与二维对称性验证}

\subsection{全局能量平衡}

全域相对总焓为
\begin{equation}
      E(t)=\Delta x\Delta y\sum_{i,j}H_{ij}(t).
      \label{eq:total-enthalpy}
\end{equation}
由四面边界热流得到单位出平面厚度的瞬时输入功率
\begin{equation}
      P_{\partial\Omega}=
      \Delta y\sum_j
      \left(q^x_{-\frac12,j}-q^x_{N_x-\frac12,j}\right)
      +\Delta x\sum_i
      \left(q^y_{i,-\frac12}-q^y_{i,N_y-\frac12}\right).
      \label{eq:boundary-power}
\end{equation}
程序用与单元更新完全相同的边界热流累计
\begin{equation}
      Q_{\partial\Omega}^{n+1}
      =Q_{\partial\Omega}^{n}+\Delta t_nP_{\partial\Omega}^{n},
      \label{eq:boundary-heat}
\end{equation}
并报告能量残差
\begin{equation}
      R_E(t)=E(t)-E(0)-Q_{\partial\Omega}(t).
      \label{eq:energy-residual}
\end{equation}
理想离散代数中 $R_E=0$；实际浮点运算只留下舍入误差。极小残差说明内部共享面热流正确抵消、
边界输入与储能更新一致，但由于二者来自同一组面热流，它不是独立解析真值，也不能单独证明
网格已足够细或导热模型完全符合实验。

\subsection{固相面积、融化比例和等效融化深度}

液相率是单元体积平均，因此连续固相面积定义为
\begin{equation}
      A_s(t)=\Delta x\Delta y\sum_{i,j}(1-f_{\ell,ij}),
      \label{eq:solid-area}
\end{equation}
而不是先用 $\fl=0.5$ 画轮廓再二值计数。相对于初始面积 $A_{s,0}=w_i h_i$，融化比例为
\begin{equation}
      F_m(t)=1-\frac{A_s(t)}{A_{s,0}}.
      \label{eq:melted-fraction}
\end{equation}
为用一个长度描述二维融化量，定义等效均匀内蚀深度 $d_{\mathrm{eq}}$：假想初始矩形四边均
向内后退相同距离 $d_{\mathrm{eq}}$，使剩余矩形与真实连续固相面积相等，亦即
\begin{equation}
      (w_i-2d_{\mathrm{eq}})(h_i-2d_{\mathrm{eq}})=A_s.
      \label{eq:equivalent-depth-definition}
\end{equation}
取物理上较小的根得到
\begin{equation}
      d_{\mathrm{eq}}
      =\frac{w_i+h_i-
      \sqrt{(w_i-h_i)^2+4A_s}}{4}.
      \label{eq:equivalent-depth}
\end{equation}
对方冰，等效边长为 $\ell_{\mathrm{eq}}=\sqrt{A_s}$，于是
$d_{\mathrm{eq}}=(w_i-\sqrt{A_s})/2$。$d_{\mathrm{eq}}$ 是面积积分得到的等效标量，并不表示
实际二维界面每一处都后退同样距离。

程序另报告水平、竖直中心线的固相率积分。以偶数行数的水平中心线为例，先平均中间两行：
\begin{equation}
      W_{\mathrm{mid}}
      =\Delta x\sum_i
      \left[1-\frac{f_{\ell,N_y/2-1,i}+f_{\ell,N_y/2,i}}{2}\right].
      \label{eq:midline-width}
\end{equation}
中心线宽度与面积等效边长一般不同；二者的差异正能反映角部比边中部更快融化的二维效应。

\subsection{反射与方形的 \texorpdfstring{$D_4$}{D4} 对称性}

若域、网格、初始冰和边界条件完全对称，则数值解应保持相应对称性。对于方形情形，正方形的
八种旋转与反射构成 $D_4$ 群。对任意场 $X$ 定义
\begin{equation}
      \epsilon_{D_4}(X)
      =\max_{g\in D_4}\lVert X-g(X)\rVert_{\infty}.
      \label{eq:d4-error}
\end{equation}
温度、焓和液相率量纲不同，程序用各自尺度归一化后取最大值：
\begin{equation}
      \epsilon_{\mathrm{sym}}=max\left[
      \epsilon_{D_4}(\fl),\,
      \frac{\epsilon_{D_4}(T)}{\max(1,|T_b-T_m|)},\,
      \frac{\epsilon_{D_4}(H)}
      {\rho[L+c_{p,w}(T_b-T_m)]}
      \right].
      \label{eq:normalized-symmetry}
\end{equation}
一般矩形算例只检查水平和竖直反射。对称误差能敏感发现数组转置、某一壁符号错误等实现问题，
但对称的错误解仍可能通过此检查，因此必须同时考察能量平衡和网格细化。

\section{默认尺度、参数与数值结果}

\subsection{为何选择该时间与空间尺度}

默认域宽高均为 $\SI{30}{\milli\metre}$，方冰边长为 $\SI{12}{\milli\metre}$，网格间距
$\Delta x=\Delta y=\SI{0.25}{\milli\metre}$。该尺度有三点作用：冰体跨 $48$ 格，界面曲率和
角部二维传热能够被清晰分辨；每侧 $\SI{9}{\milli\metre}$ 的水浴隔开冰与外壁；同时计算域
仍足够小，可用纯 NumPy 显式推进。水浴升温至 $\SI{20}{\degreeCelsius}$，使
$\SI{180}{\second}$ 内已有约 $44\%$ 初始固相面积融化，既能直观看到界面移动，又没有让方冰
完全消失。

\begin{table}[htbp]
      \centering
      \caption{默认几何、初边值与材料参数。}
      \label{tab:default-parameters}
      \begin{tabular}{lll}
            \toprule
            类别 & 参数 & 默认值 \\
            \midrule
            几何 & $L_x\times L_y$ & $\SI{30}{\milli\metre}\times\SI{30}{\milli\metre}$ \\
                 & $w_i\times h_i$ & $\SI{12}{\milli\metre}\times\SI{12}{\milli\metre}$ \\
            网格 & $N_x\times N_y$ & $120\times120$ \\
                 & $\Delta x=\Delta y$ & $\SI{0.25}{\milli\metre}$ \\
            温度 & $T_m$；$T_b$ & $\SI{0}{\degreeCelsius}$；$\SI{20}{\degreeCelsius}$ \\
            密度 & $\rho$ & $\SI{1000}{\kilogram\per\metre\cubed}$ \\
            比热 & $c_{p,w}$；$c_{p,i}$ & $4186$；$2100\ \si{\joule\per\kilogram\per\kelvin}$ \\
            导热率 & $k_w$；$k_i$ & $0.60$；$2.20\ \si{\watt\per\metre\per\kelvin}$ \\
            潜热 & $L$ & $\SI{334000}{\joule\per\kilogram}$ \\
            时间 & $t_{\mathrm{end}}$；输出间隔 & $\SI{180}{\second}$；$\SI{30}{\second}$ \\
            控制数 & $\mathrm{Fo}$ & $0.15$ \\
            \bottomrule
      \end{tabular}
\end{table}

由上述参数得到
\begin{align}
      \alpha_w&=\frac{k_w}{\rho c_{p,w}}
      =\SI{1.4333493e-7}{\metre\squared\per\second},\\
      \alpha_i&=\frac{k_i}{\rho c_{p,i}}
      =\alpha_b
      =\SI{1.0476190e-6}{\metre\squared\per\second},\\
      \Delta t_{\max}&=\SI{9.9431818e-3}{\second},\\
      \Delta t_0&=\SI{8.9488636e-3}{\second}.
      \label{eq:default-derived}
\end{align}

\subsection{默认运行结果及其解释}

默认输出时刻为 $t=(0,30,60,90,120,150,180)\ \si{\second}$，因每段末尾的短步处理，总步数
为 $20118$。$t=\SI{180}{\second}$ 的主要结果列于表~\ref{tab:default-results}。

\begin{table}[H]
      \centering
      \caption{默认算例在 $t=\SI{180}{\second}$ 的积分诊断。}
      \label{tab:default-results}
      \begin{tabular}{lll}
            \toprule
            诊断量 & 数值 & 含义 \\
            \midrule
            $A_s$ & $\SI{80.200387709}{\milli\metre\squared}$ & 连续固相率积分面积 \\
            $F_m$ & $44.3052863\%$ & 初始固相面积的融化比例 \\
            $\ell_{\mathrm{eq}}$ & $\SI{8.955466917}{\milli\metre}$ & 面积等效正方形边长 \\
            $d_{\mathrm{eq}}$ & $\SI{1.522266541}{\milli\metre}$ & 面积等效均匀内蚀深度 \\
            $W_{\mathrm{mid}}=H_{\mathrm{mid}}$ & $\SI{9.301623864}{\milli\metre}$ & 两条中心线固相积分宽度 \\
            $Q_{\partial\Omega}$ & $\SI{6095.418250235}{\joule\per\metre}$ & 四壁累计输入热量 \\
            $R_E$ & $-4.37\times10^{-11}\ \si{\joule\per\metre}$ & 全局能量残差 \\
            相对能量残差 & $7.16\times10^{-15}$ & $|R_E|/\max(1,|Q_{\partial\Omega}|)$ \\
            $\epsilon_{\mathrm{sym}}$ & 机器级零 & 默认方形条件的归一化对称误差 \\
            \bottomrule
      \end{tabular}
\end{table}

最终中心单元仍为纯冰，$f_{\ell,\mathrm{center}}=0$；单元中心温度范围为
$0\le T\le\SI{19.9981907}{\degreeCelsius}$。四壁的 $\SI{20}{\degreeCelsius}$ 温度施加在
边界面上，而非第一层单元中心，因此不应要求单元中心最大温度恰等于 $\SI{20}{\degreeCelsius}$。
此外，$W_{\mathrm{mid}}$ 大于 $\ell_{\mathrm{eq}}$，说明角部融化快于边中部；二维冰块并未简单保持
为均匀缩小的正方形。

为检查空间细化趋势，保持物理尺寸和时间不变，在 $t=\SI{60}{\second}$ 比较三个方形网格：
\begin{table}[H]
      \centering
      \caption{$t=\SI{60}{\second}$ 的网格细化趋势。}
      \label{tab:grid-refinement}
      \begin{tabular}{ccc}
            \toprule
            $N_x=N_y$ & $\Delta x$（\si{\milli\metre}） & $d_{\mathrm{eq}}$（\si{\milli\metre}） \\
            \midrule
            80  & 0.3750 & 0.817989 \\
            120 & 0.2500 & 0.810287 \\
            160 & 0.1875 & 0.806112 \\
            \bottomrule
      \end{tabular}
\end{table}
后两级之差小于前两级之差，表明该积分量随网格细化呈稳定趋势。三组结果尚不足以证明理论
收敛阶、进入渐近收敛区或收敛到解析真值，故本文不作此类过度推断。

\section{运行、输出与结果复核}

在项目根目录执行
\begin{lstlisting}[style=iceflowcode]
python iceflow2d/examples/stefan_melting_2d.py
\end{lstlisting}
即可使用默认参数运行。命令行选项允许修改域尺寸、网格数、冰尺寸、温度、材料参数、
Fourier 数、时间步和输出位置；显式指定的时间步仍须通过稳定性检查。默认输出目录为
\path{outputs/iceflow2d/stefan_melting_2d}，包含：
\begin{itemize}
      \item \code{history.csv}：七个输出时刻的面积、融化量、能量和对称性诊断；
      \item \code{fields.npz}：坐标、时间及形状为 $(7,120,120)$ 的 $T,\fl,H$ 场；
      \item \code{metadata.json}：输入配置、派生尺度、模型假设和最终诊断；
      \item \code{stefan\_melting\_2d.png}：四个时刻的液相率、最终温度和融化历史汇总图。
\end{itemize}
图中的 $\fl=0.5$ 等值线只服务于可视化；求解器没有把它当作显式追踪界面，固相面积也不由该
二值轮廓计算。复核结果时，建议依次检查：配置记录是否与预期一致；$0\le\fl\le1$；融化量和
温度场是否保持几何对称；式~\eqref{eq:energy-residual} 是否接近浮点舍入量级；网格细化后积分
诊断是否趋于稳定。

\section{方法边界与结论}

本算例把融化问题转化为一个清晰的守恒循环：体积焓保存能量，分段闭合将焓解释为温度和液相
率，有限体积面热流负责在相邻单元间交换能量，显式更新则推进时间。由于潜热包含在 $H$ 中，
冰--水界面可在固定网格上自然穿过单元，无需为界面位置建立额外运动方程。

该实现适合作为二维相变热模块的基础验证，但结论受其假设限制。它没有流体运动，因而不能表现
自然对流对融化形状和速率的影响；采用等密度，因而不能表现融化收缩或结冰膨胀；恒温外壁和
固定 $k_w$ 边界半格也只对应当前热水浴场景。若将来与 LBM、真实 $\rho_i\ne\rho_w$、自由液面
或移动刚体耦合，应在保留本文共享面通量与焓守恒结构的同时，另行加入显热对流、相变几何更新、
质量--体积转换和流体边界重建。本文的默认结果首先证明独立二维热焓实现具有内部能量守恒、
预期对称性和合理的空间细化趋势，而不把这些自洽检查误称为完整实验验证或解析解验证。

\end{document}
